\documentclass{article}

\usepackage[final]{neurips_2026}
\makeatletter
\renewcommand{\@noticestring}{Workshop on Who Verifies the Agents? Toward Reliable Agent Development, NeurIPS 2026.}
\makeatother

\usepackage[utf8]{inputenc}
\usepackage[T1]{fontenc}
\usepackage{hyperref}
\usepackage{url}
\usepackage{booktabs}
\usepackage{amsfonts}
\usepackage{amsmath}
\usepackage{amssymb}
\usepackage{microtype}

\title{Breaking the Composition Cliff: Mutation-Guided Specification Synthesis for Multi-Function Programs\\[0.5em] \large \mdseries --- Abstract ---}

\author{%
  Ishraq Tanvir \\
  FOF Business Point \\
  \texttt{ishraq.tanvir@fofbd.com} \\
}

\begin{document}

\maketitle

\begin{abstract}
Large language models (LLMs) can generate code, but synthesizing formal specifications for this code remains challenging. While previous approaches report high verification rates on individual functions, real-world programs are composed of modular interprocedural call graphs. We evaluate six frontier language models on synthesizing formal Dafny specifications for 300 multi-function programs with call-graphs of sizes 2--5 (\textsc{DafnyComp-300}), revealing two key shortcomings. First, we find that models struggle with compositional specifications: while models reach 91.7\%--100\% specification completeness on individual functions with repair, zero-shot specification success collapses from up to 95.0\% on single functions down to 14.0\%--41.3\% on multi-function call graphs. Second, the traditional repair loop, which relies on compiler error messages, leads models to learn vacuous specifications, e.g., trivial inequalities such as \texttt{ensures res >= 0}. While these specifications appear logically sound, they are overly weak and cannot reject incorrect implementations. To address this, we present \textsc{Mutation-CEGIS}, an approach that uses AST code mutants to test specifications for vacuity. When a candidate specification mistakenly admits a buggy mutant, we extract the code change as a concrete counterexample for the model to refine its specification. On single-function benchmarks (\textsc{CloverBench-60}), \textsc{Mutation-CEGIS} boosts specification completeness to 91.7\%--100.0\% while reducing repair turns by up to 36.8\% and saving up to 24.9\% of total tokens. On compositional call graphs (\textsc{DafnyComp-300}), \textsc{Mutation-CEGIS} systematically detects vacuous bounding contracts across model families and, when synthesis succeeds, drives models toward tighter first-order relational specifications ($\forall, \exists$). Finally, we find that for deeper call-graphs, exact specification synthesis results in SMT-solver timeouts, demonstrating that future verification agents will benefit from synthesizing auxiliary lemmas in addition to specifications.
\end{abstract}

\vspace{2em}
\noindent\textbf{Keywords:} Formal Verification, Program Synthesis, CEGIS, Mutation Testing, Deductive Verification, Code LLMs, AI Agents, Agent Verification, Dafny, SMT Solvers, Z3, Hoare Logic, Invariant Synthesis, Software Reliability.

\vspace{1em}
\noindent\textbf{TL;DR:} We uncover the Composition Cliff in LLM formal verification and present \textsc{Mutation-CEGIS} using AST code mutants to eliminate vacuous contracts and synthesize hardened first-order specifications across multi-function programs.

\vspace{1em}
\noindent\textbf{Code \& Benchmarks:} \url{https://github.com/devishraq/composition-cliff}

\end{document}
